\section{From Stateful Agents to Transactional High-Capability Reasoning}
\label{sec:transactional}

RaS treats a high-capability reasoning episode as a bounded computation over externally persisted engineering state:
\begin{equation}
R_t + U_t
\rightarrow
\text{reasoning episode}
\rightarrow
\text{validated state transition}
\rightarrow
R_{t+1}.
\end{equation}
After $R_{t+1}$ is accepted, the reasoning process may be terminated and replaced.

The term \emph{transactional} is retained only as a lifecycle analogy. It does not imply ACID semantics, serialisability, transactional memory, deterministic replay, or database atomicity. Recent work on semantic isolation for durable AI workflows uses transaction terminology for a different problem: maintaining compatible semantic resources across paused and resumed workflows \citep{mozafari2026semantic}. RaS therefore uses “reasoning transaction” narrowly to mean a bounded reasoning episode over externally persisted authoritative state whose process identity need not survive after an accepted transition.

\subsection{Disposable reasoning is recomputation-tolerant, not computation-free}

A reasoning process is \emph{disposable with respect to programme continuity} if replacing it after an accepted transition does not materially impair the next correct programme continuation under a prespecified behavioural comparison, given matched durable state, task, model capability/configuration, tool interface, and execution environment.

This definition explicitly allows recomputation. A fresh reasoner may re-read files, rediscover dependencies, or re-establish architectural understanding. If that work is large, disposal may be technically feasible but economically inferior. “Disposable” therefore concerns lifecycle dependence, not elimination of reasoning cost.

Figure~\ref{fig:lifecycle} contrasts the persistent-history control with the forced-reset lifecycle.

\begin{figure}[htbp]
\centering
\begin{tikzpicture}[
node distance=5mm,
box/.style={draw,rounded corners,align=center,minimum width=3.35cm,minimum height=7mm,font=\small},
arrow/.style={-{Latex[length=1.8mm]},thick}
]
\node[font=\bfseries] at (-3.25,1.0) {A: persistent-history control};
\node[box] (p1) at (-3.25,0) {$R_0+U_1$\\reasoner session $A$};
\node[box,below=of p1] (p2) {accept $R_1$; rematerialise\\canonical workspace};
\node[box,below=of p2] (p3) {$R_1+U_2$ + predecessor\\reasoning history};
\draw[arrow] (p1)--(p2);
\draw[arrow] (p2)--(p3);

\node[font=\bfseries] at (3.25,1.0) {B: forced reset};
\node[box] (r1) at (3.25,0) {$R_0+U_1$\\fresh session $B_1$};
\node[box,below=of r1] (reset1) {accept $R_1$; destroy $B_1$;\\rematerialise same workspace};
\node[box,below=of reset1] (r2) {$R_1+U_2$\\fresh session $B_2$};
\draw[arrow] (r1)--(reset1);
\draw[arrow] (reset1)--(r2);

\node[below=12mm of p3,align=center,font=\small] (note) {Model/configuration, tools, environment, task, accepted repository state, and verifier are matched.\\The intended treatment difference is predecessor reasoning-session history.};
\end{tikzpicture}

\caption{Conceptual lifecycle comparison. The experimental treatment is predecessor reasoning-session history at an accepted engineering boundary; workspace, model, tools, task, and accepted repository state must otherwise be matched.}
\label{fig:lifecycle}
\end{figure}

Checkpoint/restart and stateless-service systems are important precedents rather than novelty claims. Durable Functions, for example, demonstrates reliable workflow execution with persisted state and reconstructable compute \citep{burckhardt2021durable}. Git already provides checkpoint-like accepted repository states. The RaS-specific question is semantic: whether those accepted software artefacts contain enough information that a fresh high-capability reasoner can continue a \emph{dependent engineering programme} without predecessor conversational/model-native history.

Potential properties such as restartability, horizontal scheduling, model routing, and lower session affinity are architectural implications only. Model routing is established prior work \citep{ong2024routellm}; cross-model continuity is not tested in P0 and remains future work.

The architecture also exposes failure classes. A reset run can fail because reconstruction selected the wrong evidence; because the reasoner misinterpreted correct evidence; because the proposal omitted a constraint; because execution failed; because the environment changed; or because validation was invalid. These mechanisms must not be collapsed into one causal explanation.

The central experimental requirement is therefore stronger than “open a new chat”. The audit requires canonical workspace rematerialisation and invariant system instructions, model configuration, task specification, tool interface, dependency state, and environment across conditions. If hidden provider/session memory or uncontrolled tool state differs between A and B, the treatment is not identifiable.
